\section{Discussion}
\label{sec:discussion}

\subsection{Legal Theory: CSS as a Litigation Tool}
\label{sec:legal}

The Census Stability Score has a potential evidentiary application in
post-\emph{Rucho} partisan gerrymandering litigation.

\paragraph{The constitutional argument.}
Under \emph{Rucho v.\ Common Cause} (2019), federal courts cannot adjudicate
partisan gerrymandering claims.
But state courts, applying their own constitutions, retain this jurisdiction.
The key factual question in state-court litigation is: \emph{does the enacted
map deviate from the natural map, and if so, can the deviation be explained
by legitimate geographic factors?}

A GeoSection map with CSS $\geq 0.90$ would provide a strong geography-first
record for this question.
It would show that the same procedure selected the same structure across three
independent decennial censuses.  Any enacted map that differs from that
structure would then need to be explained against a concrete benchmark, but CSS
alone would not decide the legal question.

\begin{conjecture}[CSS as Geographic Evidence]
\label{prop:burden-shift}
If a state has a GeoSection map with CSS$_{3\text{yr}} \geq 0.90$, then
an enacted congressional map that differs from the GeoSection map should be
compared against that stable benchmark before the deviation is attributed to
geography alone.\footnote{This is a proposed evidentiary use of CSS, not a
claim about the governing legal burden in any jurisdiction.  The full
three-census package remains future evidence.}
\end{conjecture}

The strength of Conjecture~\ref{prop:burden-shift} increases with:
(a) the number of census years confirming the same outcome,
(b) the magnitude of the proportionality gap in the enacted map relative
to the GeoSection map, and
(c) the stability of the surrounding states (if all neighbouring states
have CSS $\geq 0.90$, the geographic argument is regional, not idiosyncratic).

\paragraph{The historical-vintage argument.}
A map with CSS $\geq 0.90$ across 2000, 2010, and 2020 existed before the
latest census-vintage political data used in the current redistricting cycle.
In 2000, the specific partisan geography that characterises today's Republican
dominance in exurban and rural areas was less pronounced; the urban-rural
divide was real but less extreme.
If GeoSection produced the same 5\,D/9\,R outcome in 2000 as in 2020 for a
state like Wisconsin, that would be evidence that the outcome is not contingent
on the current partisan distribution.
It predates the 2016 political realignment.
It predates the post-2010 Tea Party Republican gains in suburban areas.
It would be geography-first evidence, subject to the usual caveats about
election data, enacted-plan comparison, and state law.

\paragraph{The Lorenz-stable map as a statute design tool.}
A proposed Districting Integrity Act could incorporate CSS as a procedural
threshold:
\begin{enumerate}
\item Any proposed congressional map must be accompanied by its CSS score
  computed against the GeoSection-optimal map.
\item Maps with CSS $\geq 0.90$ may be eligible for administrative
  certification under a statute that chooses to create that pathway.
\item Maps with CSS $0.70$--$0.90$ require a majority-vote legislative approval
  with public disclosure of the CSS components.
\item Maps with CSS $< 0.70$ trigger a heightened judicial review process,
  because the proposed map deviates substantially from the state's durable
  geographic structure.
\end{enumerate}

This framework respects democratic self-governance: legislatures retain the
ability to adopt non-GeoSection maps, but they must publicly acknowledge the
degree to which their map departs from the natural geographic partition.

\paragraph{States with high CSS are the clearest review targets.}
When a high-CSS state has an enacted map that differs from GeoSection,
the discrepancy is harder to explain by geography alone.
Wisconsin (CSS $\approx 0.90$, predicted) has a 1:7 ratio anchored by
Milwaukee's fixed municipal boundary and Lake Michigan shoreline.
If Wisconsin's enacted map packed Milwaukee into a single district in a
way that produced a different city-versus-suburbs configuration, the
discrepancy would require a more specific geographic explanation: the
algorithm, across three censuses, would have consistently identified
Milwaukee's compact geography as the natural first-level peel.
Note: Iowa is a low-CSS state (confirmed ratio shift from $1:3$ in 2000 to
$2:2$ in 2010; see \S\ref{sec:iowa-case-study}) and is therefore a weaker
litigation target under this framework despite its recent near-uniform
agricultural geography.

\subsection{Null Hypothesis: Expected Stability Under Random Ratio Assignment}
\label{sec:null-hypothesis}

The null hypothesis for CSS stability is a randomly drawn partition changing
its natural ratio with probability $p_{\text{null}}$.
For a state with $k$ districts trying all $\lfloor k/2 \rfloor$ candidate
ratios, the probability of selecting the same ratio twice by chance (assuming
a uniform distribution over candidate ratios) is $2/k$.
Across our 30-state sample, weighting each state by its seat count, the
expected random stability rate is approximately 20\,\%.

Our observed 67\,\% stability rate is $3.4\times$ the null expectation,
providing strong evidence that census-stability reflects geographic
determinism rather than chance.
Put differently: if GeoSection selected its first-level ratio uniformly at
random from the candidate set in each census year, we would expect
approximately 6 of 30 states to show spurious stability.
We observe 20 stable states, a surplus of 14 over the random baseline.
The binomial $p$-value for observing $\geq 20$ stable states under the 20\,\%
null is $p < 0.001$ ($\binom{30}{k} 0.2^k 0.8^{30-k}$ for $k \geq 20$).

This framing suggests that the 67\,\% finding is not merely an artefact of a
lenient stability definition.
Even using the conservative 5\,\% threshold ($\tau = 0.05$), the observed
rate is far above the null, and the lower bound of the Wilson 95\,\%
confidence interval (48\,\%) remains well above the 20\,\% null.

\subsection{Population Dynamics and Geographic Reorganisation}
\label{sec:pop-dynamics}

States with low CSS are not failed cases for GeoSection.
They are states where the population-geography relationship has genuinely
reorganised between census cycles.

\paragraph{The Sun Belt restructuring hypothesis.}
Arizona's Phoenix metro grew from approximately 2.8M people in 2000 to
4.9M in 2020---a 75\,\% increase in twenty years.
This growth extended the Phoenix metro boundary into the far East Valley,
the West Valley, and North Scottsdale, creating new geographic regions that
are contiguous with Phoenix's compact urban core but have distinct
socioeconomic and demographic characteristics.
For GeoSection, this means the isoperimetric boundary of ``the Phoenix urban
core'' has expanded: in 2000, the natural partition may have peeled a compact
2.8M-person metro from a 5.1M total state population ($p^* \approx 0.13$, $1:7$);
in 2020, the metro is 4.9M of 7.2M total ($p^* \approx 0.11$, $1:8$).
The ratio shift from $1:7$ to $1:8$ is not an algorithm failure.
It is a signal that Arizona's human geography has reorganised.

\paragraph{Declining states: stability through contraction.}
Ironically, states with declining populations may exhibit high CSS because
geographic concentration \emph{decreases} over time.
West Virginia, Mississippi, and Alabama have seen substantial rural
depopulation, with population concentrating in smaller areas.
This tends to sharpen the geographic features that GeoSection detects,
potentially increasing the isoperimetric gap between ratios and making the
natural ratio more robust.
We predict these states will have high CSS despite (or because of)
demographic decline.

\paragraph{The Lorenz drift proxy as an early warning.}
States where $p^*$ shifts between census years are states where the
population-geography Lorenz curve has restructured.
Monitoring $\Delta p^*$ between the ACS 5-year estimates (available annually)
and the decennial census provides an early warning of potential ratio shifts.
A state where $\Delta p^*$ exceeds $0.03$ between two ACS estimates should
flag for updated GeoSection analysis, potentially before the decennial
redistricting trigger.

\subsection{Limitations}

\paragraph{Tract alignment across census years.}
The Jaccard similarity computation requires mapping district assignments from
2000 and 2010 tract boundaries to 2020 boundaries.
Our spatial intersection method (population-weighted allocation to 2020 tracts)
introduces approximation error in areas where the 2000 tract network differs
significantly from 2020.
Urban cores, where tracts are frequently split, have the highest
approximation error.
This means Jaccard similarity is slightly underestimated in urban areas,
potentially underestimating CSS for states with large cities.

\paragraph{Population-only perturbation requires common graph.}
The Type I/Type II decomposition requires running GeoSection on the 2020
graph with 2000 populations.
This is an approximation: 2000 tract population totals must be allocated to
2020 tracts by spatial intersection, which introduces the same approximation
error as the Jaccard computation.
For the stability \emph{difference} (full cross-year minus population-only
perturbation), the approximation errors partially cancel; but they do not
cancel exactly.

\paragraph{Seed count for large states.}
California ($k=52$), Texas ($k=38$), and Florida ($k=28$) have uncertain
natural ratios at 50 seeds per ratio.
Cross-census stability for these large states requires higher seed counts
(200+ per ratio) to ensure convergence.
We report their 2020 natural ratios as provisional and exclude them from the
CSS computation until higher-confidence estimates are available.

\paragraph{Package-complete CSS not yet available.}
The three-census CSS requires the year-complete package and a common cross-year
comparison format.
Two-census CSS covers the largest twenty-year span (2000--2020), but it is not
yet a substitute for the three-census table because it omits the intermediate
redistricting cycle.
Three-census CSS will strengthen the claim by adding an intermediate
data point that lies \emph{between} the bookend years.

\paragraph{Congressional apportionment changes the comparison baseline.}
Fourteen states changed seat count between 2000 and 2020.
Our normalised ratio comparison addresses this, but the comparison is
intrinsically less direct than for states with fixed seat counts.
A state like Texas ($k=30$ in 2000, $k=38$ in 2020) has a ratio comparison
that conflates algorithm stability with apportionment-change effects.
The CSS for seat-count-changing states should be interpreted with this caveat.

\subsection{The 2D Stability Picture: (seed, census)}

Combining B.7's seed stability results with T.8's census stability evidence and predictions,
we can sketch a $2\times 2$ stability matrix:

\begin{table}[h]
\centering
\caption{$2\times 2$ stability matrix: seed stability (B.7) $\times$ census
  stability (T.8 evidence and predictions). Representative states in each cell.}
\label{tab:2d-stability}
\begin{tabular}{l l l}
\toprule
                  & Low census stability & High census stability \\
\midrule
High seed stability & AZ, NV, CO, IA      & KS, WI, IL, PA \\
Low seed stability  & TX, FL, CA         & (none expected) \\
\bottomrule
\end{tabular}
\end{table}

States in the top-right cell (high seed stability, high census stability)
have the strongest geographic determinism record and are the easiest to defend
as geography-first outputs.
States in the bottom-right cell are expected to be rare:
if the geographic structure is durable across censuses, it should also be
easily detectable by the algorithm (high seed convergence).
States in the bottom-left cell (TX, FL, CA) are algorithmically difficult:
their geographic structure is complex enough that multiple equally-good
solutions exist within a single census year, \emph{and} the optimal solution
shifts across census years as population redistributes.
These states require the most procedural justification for any particular map.

\subsection{Predictive Model for CSS}

The theoretical framework suggests a predictive regression:
\begin{align*}
  \mathrm{CSS}(s) \sim f(&\text{population growth},\;
                         \text{urban boundary growth},\\
                       &\text{seat-count change},\;
                         p^*_{2020}).
\end{align*}
The expected signs are negative for population growth, urban boundary growth,
and seat-count change.  The sign of $p^*_{2020}$ depends on whether the ratio
is in the metropolitan-fixity region.

This regression, fit on the package-complete three-census CSS dataset once available,
would allow CSS prediction for future census cycles before running the full
50-state sweep.
